\documentclass{amsart}
\usepackage{amsmath,amssymb,amsthm,hyperref}
\theoremstyle{plain}
\newtheorem{theorem}{Theorem}
\newtheorem{proposition}{Proposition}
\newtheorem{lemma}{Lemma}
\theoremstyle{remark}
\newtheorem{remark}{Remark}
\newtheorem*{statusnote}{Status}
\DeclareMathOperator{\rad}{rad}
\begin{document}

\title{On the equation $x(x+1)=y!$}
\maketitle

\begin{abstract}
We study the Diophantine equation $x(x+1)=y!$ in positive integers.
The only solutions presently known are $(x,y)=(1,2)$ and $(x,y)=(2,3)$,
and we verify computationally that there are no others with $y\le 2000$.
We prove the genuine facts that can be established rigorously: the
reformulation of the equation as the statement that $4\,y!+1$ is a perfect
square, the basic structure of the two coprime factors, and the fact that
the equation has only finitely many solutions (an ineffective theorem of
Berend--Osgood).  We also explain precisely why the usual elementary
obstructions (congruences and prime--valuation parity) \emph{cannot} settle
the problem.  We emphasize honestly that \textbf{no proof is known} that
$(1,2)$ and $(2,3)$ are the only solutions; this is an open problem of
Brocard--Ramanujan type.
\end{abstract}

\section{Statement and summary}

We consider
\begin{equation}\label{eq:main}
   x(x+1)=y!,\qquad x,y\in\mathbb{Z}_{\ge 1},
\end{equation}
where $y!=1\cdot 2\cdots y$.

\begin{statusnote}
The question of \emph{determining all} solutions of \eqref{eq:main} is, to
the best of current knowledge, \textbf{open}.  The only solutions known are
\[
   (x,y)=(1,2)\quad\bigl(1\cdot 2=2!\bigr),\qquad
   (x,y)=(2,3)\quad\bigl(2\cdot 3=3!\bigr),
\]
and exhaustive computation (Section~\ref{sec:comput}) shows there is no
further solution with $y\le 2000$.  It is a theorem (Berend--Osgood,
Section~\ref{sec:finite}) that \eqref{eq:main} has only finitely many
solutions, but that theorem is \emph{ineffective}: it does not allow one to
conclude that the list above is complete.  We do not claim a complete
proof, and we explain in Section~\ref{sec:obstruction} why the elementary
techniques that suffice for many factorial Diophantine equations are
provably insufficient here.
\end{statusnote}

The remainder of the paper records exactly what can be proved rigorously.

\section{Reformulation and elementary structure}

\begin{proposition}\label{prop:square}
A positive integer $y$ admits a solution $x$ of \eqref{eq:main} if and only
if $4\,y!+1$ is a perfect square.  In that case the solution is unique and
given by $x=\tfrac{1}{2}\bigl(\sqrt{4\,y!+1}-1\bigr)$.
\end{proposition}

\begin{proof}
Equation \eqref{eq:main} is the quadratic $x^2+x-y!=0$ in $x$, whose
discriminant is $1+4\,y!$.  A positive integer root exists iff
$\sqrt{1+4\,y!}\in\mathbb{Z}$ (equivalently $1+4\,y!$ is a perfect square),
because in that case
\[
   x=\frac{-1+\sqrt{1+4\,y!}}{2}
\]
is the unique positive root; the value is an integer since $1+4\,y!$ is
odd, so $\sqrt{1+4\,y!}$ is odd and $\sqrt{1+4\,y!}-1$ is even.  Conversely
if $x$ solves \eqref{eq:main} then $(2x+1)^2=4x^2+4x+1=4\,y!+1$, a perfect
square.  Uniqueness of $x$ for each $y$ is clear since $x\mapsto x(x+1)$ is
strictly increasing on $\mathbb{Z}_{\ge 1}$.
\end{proof}

Thus \eqref{eq:main} is equivalent to the requirement that $4\,y!+1$ be a
square; with $m=2x+1$ it reads
\begin{equation}\label{eq:m}
   m^2=4\,y!+1,\qquad m\ \text{odd}.
\end{equation}
This places \eqref{eq:main} squarely in the Brocard--Ramanujan family
``$c\cdot n!+1=\square$''; the classical Brocard problem is the case
$n!+1=\square$, which is itself famously open.

\begin{lemma}\label{lem:coprime}
If $x(x+1)=y!$ then $\gcd(x,x+1)=1$, so $x$ and $x+1$ are coprime divisors
of $y!$ whose product is $y!$.  Consequently, for every prime $p\le y$, the
entire $p$--part $p^{\,v_p(y!)}$ of $y!$ divides exactly one of $x$, $x+1$.
\end{lemma}

\begin{proof}
Consecutive integers are coprime: any common divisor divides their
difference $1$.  Since $x(x+1)=y!$ and $\gcd(x,x+1)=1$, each maximal prime
power $p^{v_p(y!)}\,\|\,y!$ must divide a single one of the two coprime
factors (it cannot split across them, for then $p$ would divide both).
\end{proof}

\section{Why the standard elementary obstructions fail}\label{sec:obstruction}

It is worth recording, as a rigorous negative result, that the two
techniques which most often resolve factorial equations are powerless here.

\paragraph{Congruence (modular) obstructions do not exist.}
For any fixed modulus $N$ and all $y\ge N$ we have $N\mid y!$, hence
$4\,y!+1\equiv 1\pmod{N}$.  Since $1$ is a quadratic residue modulo every
$N$, equation \eqref{eq:m} has no local obstruction at any modulus for
large $y$.  In particular one cannot rule out solutions by congruences.

\paragraph{Valuation--parity obstructions do not exist.}
A common idea is: if some prime $p$ has $v_p(y!)$ \emph{odd}, perhaps a
factor is forced to be a non-square and a contradiction follows.  This
fails because Lemma~\ref{lem:coprime} only forces $p^{v_p(y!)}$ to divide
one of $x,x+1$; neither $x$ nor $x+1$ is required to be a perfect square
(only their \emph{product} $y!$ is constrained, and $y!$ is generally not a
square).  Indeed in the genuine solution $2\cdot 3=3!$ we have $v_2(3!)=1$
and $v_3(3!)=1$ both odd, yet a solution exists.  Hence no parity argument
on $v_p(y!)$ can preclude solutions.

These two observations explain why \eqref{eq:main}, despite its innocent
appearance, resists elementary attack and why a determination of all
solutions requires (currently unavailable) deeper input.

\section{Finiteness of the solution set}\label{sec:finite}

\begin{theorem}[Berend--Osgood]\label{thm:finite}
Let $P\in\mathbb{Z}[t]$ have degree $\ge 2$.  Then the equation $P(x)=y!$
has only finitely many solutions $(x,y)\in\mathbb{Z}_{\ge 1}^2$.  In
particular, $x(x+1)=y!$ has only finitely many solutions.
\end{theorem}

This is the main result of D.~Berend and C.~Osgood, \emph{On the equation
$P(x)=n!$ and a question of Erd\H{o}s}, J.\ Number Theory \textbf{42}
(1992), 189--193.  We recall its meaning and limitation rather than
reproducing its proof.

\begin{remark}[The theorem is ineffective]
The Berend--Osgood argument shows the set of solutions is finite but does
\emph{not} produce an explicit bound on $y$.  It therefore does \emph{not}
imply that $(1,2)$ and $(2,3)$ exhaust the solutions: it leaves open the
logical possibility of a further (necessarily enormous) solution.  This is
exactly why the determination problem remains open even though finiteness
is a theorem.
\end{remark}

\section{Computational verification}\label{sec:comput}

\begin{proposition}\label{prop:comp}
For $1\le y\le 2000$ the only $y$ for which $4\,y!+1$ is a perfect square
are $y=2$ and $y=3$.  Hence the only solutions of \eqref{eq:main} with
$y\le 2000$ are $(1,2)$ and $(2,3)$.
\end{proposition}

\begin{proof}[Verification]
By Proposition~\ref{prop:square} it suffices to test, for each
$y\in\{1,\dots,2000\}$, whether $4\,y!+1$ is a perfect square; this is
decided exactly in integer arithmetic by computing
$s=\lfloor\sqrt{4\,y!+1}\rfloor$ and checking $s^2=4\,y!+1$.  Maintaining
$y!$ incrementally, the check is exact (no floating point), and it returns
``square'' precisely for $y=2$ ($4\cdot 2+1=9=3^2$, giving $x=1$) and
$y=3$ ($4\cdot 6+1=25=5^2$, giving $x=2$).  All other $y$ in the range fail.
\end{proof}

\section{Conclusion}

Collecting the rigorous statements:
\begin{itemize}
\item By Proposition~\ref{prop:square}, \eqref{eq:main} is equivalent to
``$4\,y!+1$ is a perfect square''.
\item By Theorem~\ref{thm:finite}, \eqref{eq:main} has only finitely many
solutions; but this is ineffective (Remark after it).
\item By Proposition~\ref{prop:comp}, the only solutions with $y\le 2000$
are $(1,2)$ and $(2,3)$.
\item By Section~\ref{sec:obstruction}, neither congruence nor
valuation--parity arguments can settle the problem.
\end{itemize}
The honest conclusion is that the only \emph{known} solutions are
$(x,y)=(1,2)$ and $(x,y)=(2,3)$, that there is strong computational and
heuristic evidence (and a finiteness theorem) supporting the belief that
these are the only ones, but that a \emph{complete} determination is an
\textbf{open problem}.
\end{document}
